\documentclass[runningheads]{llncs}

% Encodage / langue (pdfLaTeX)
\usepackage[T1]{fontenc}
\usepackage[utf8]{inputenc}
\usepackage[french]{babel}

% Maths / figures / liens
\usepackage{amsmath,amssymb}
\usepackage{graphicx}
\usepackage{xcolor}
\usepackage{hyperref}

\usepackage[authoryear]{natbib}


\title{Community detection in geometric graphs: Recent advances until 2026}
\titlerunning{Community detection in geometric graphs} 
\author{Raphaël Sala}
\authorrunning{R. Sala}
\institute{Affiliation à renseigner\\ \email{raphael.sala@utoulouse.fr}}

\begin{document}
\maketitle

\begin{abstract}
\textcolor{red}{TODO}
\keywords{Community detection, geometric graphs, random geometric graphs, exact/partial recovery, spectral methods}
\end{abstract}

\section{Introduction}\

\paragraph{General idea}

Community detection consists to find clusters in a graph where nodes share more connections with nodes of the same cluster
than with nodes of other clusters. This large topic is detailled in general by \citep{fortunato2010community}. In this 
report we focus on random geometric graphs (RGG). The specificity to add spatial information allows us to consider
local information based on real or abstract proximity.  % The probability of connection between two nodes depends on their distance.
In social networks, two individuals are more susceptible to be friends if they are geographically near
or if they have similar terms of interests, age, $\dots$
For more concrete examples of RGG can be find for wireless networks from \citep{gomez2015case}, biology \citep{higham2008fitting}
and in particular an application of detection of communities \citep{sankararaman2020comhapdet}
\textcolor{red}{ word embedding, physic, ...}.

\paragraph{Why geometric} The main model for community detection is the Stochastic Block Model (SBM)
\citep{holland1983stochastic}. Each node carries a hidden community label
and edge probabilities depend on whether two nodes belong to the same community or not.
This model comes from the Erdos-Renyi graph, where each node connect with the same probability
to others nodes \citep{erdHos1960evolution}. This model is nowadays well understood,
see \citep{abbe2018community}. Nonetheless, it suffers to some drawbacks because its 
unrealistic to think that every nodes can connected to others with the same probability rate.
At the same time than Erdos-Renyi, the RGG was introduced by \citep{gilbert1961random}.
Add geometric information allows us to mimic some emergent behaviour from real network,
like sparsity \citep{Del_Genio_2011}, the "small world" property \citep{watts1998collective}
or even transitivity "friend of a friend is a friend" \textcolor{red}{Ref ?}.

\paragraph{Context of the problem}
RGG and its properties are well studied, see \citep{penrose2003random}. 
When we add community structure, we are in a different setting.
We observe the graph, and sometimes also the node locations,
while the community labels remain hidden. 
Naturally, we want to know if we can recover the latent labels, what is the best accuracy we can achieve,
and which algorithms can reach it.

\textcolor{red}{TODO
Définition et interprétation, récupération à permutation près
\begin{itemize}
    \item Distinguabilité \citep{sankararaman2018community}
    \item Récupération partielle (weak recovery)
    \item Récupération presque exacte (almost exact recovery)
    \item Récupération exacte (exact recovery) \citep{gaudio2024exact}
\end{itemize}
}
Thus, in this report, community detection is not only understood as
producing a convenient partition of the observed graph, but rather as
recovering latent community labels under a probabilistic geometric model,
up to permutation.


\paragraph{What we don't do}
Precedently we enhance properties RGG share with real networks due
to the geometric information. However,
we can get for instance transitivity, and surely more, thanks to attribute based models
\citep{Li_Sha_Huang_Zhang_2018} or 
hyper-graph \citep{Ruggeri_2023} and even with its geometric variant
\citep{Kov_cs_2025}.
We will not consider these models to focus on the geometric aspect.
Moreover, we presented problems we are interested about. From them, one
is \textit{Distinguability} to not confond with the
\textit{Detection Problem} we don't treat.
This last one ask generally 
condition on parameters, defining the RGG, to be distinguable to 
an Erdos-Renyi graph. Hence, this hasn't link to community.
Refer to 
\citep{bubeck2016testing}
\citep{liu2022testing}
\citep{mao2026randomgeometricgraphssmooth}.
Finally, at the beginning of the literature,
some works were interested to recover the latent position of nodes in the geometric space,
and consider it as a Community detection problem, \citep{valdivia2019latent}
\citep{eldan2022community}.
This is a different problem than the one we are interested in,
and it is more difficult because we want to recover the position (contnue label)
of each node rather than community (discrete label).




\section{\textcolor{red}{Modèle}}

\begin{itemize}
    \item pas de boucle, pas dirige, pas de communautés chevauchantes, etc...
    \item espace latent (ou pas) \citep{decastro2020adaptiveestimationnonparametricgeometric}
    \item sphère, espace euclidien ?
    \item régime : dense, logarithmique, sparse
    \item hard-RGG, soft-RGG \citep{duchemin2023random}
    \item motivation : bien que les modèles ne soient pas les mêmes, ils n'ont pas pour but
    de comprendre tous ceux possible mais qu'à travers chaque classe voir les mécanismes. 
    Exemplse : passer du local au global puis raffinement, les quantités utilisées, ... 
\end{itemize}
\textcolor{orange}{Le but est de faire des notaitons claires, concises et générales pour pouvori expliquer tous les modèles ffff
sans avoir à se perdre dans la notation.}

\section{\textcolor{red}{Seuil théorique information}}

\textcolor{orange}{Faire le rapport dans l'ordre chronologique je pense}

\begin{itemize}
    \item \citep{galhotra2018geometric} (Geometric Block Model, GBM)
    \item \citep{sankararaman2018community} \& \citep{abbe2021community} (Geometric SBM)
    \item \citep{chien2020active}
    \item \citep{galhotra2020connectivityrandomannulusgraphs} (GBM)
    \item \citep{galhotra2023communityrecoverygeometricblock} (GBM)
    \item \citep{avrachenkov2024community} (Geometric Kernel Block Model, GKBM)
    \item \citep{gaudio2024exactcommunityrecoverygeometric} (Geometric SBM)
    \item \citep{gaudio2024exact} (Geometric Hidden Community Model, GHCM)
    \item \citep{gaudio2025sharpexactrecoverythreshold} (GHCM)
    \item \citep{gaudio2026exact} (GHCM)
\end{itemize}

\subsection{\textcolor{red}{Geometric Block Model (GBM)}}
\citep{galhotra2018geometric} 
\citep{galhotra2020connectivityrandomannulusgraphs}
\citep{galhotra2023communityrecoverygeometricblock}

\citep{chien2020active}


\subsection{\textcolor{red}{Geometric Stochastic Block Model (GSBM)}}
\citep{sankararaman2018community} \citep{abbe2021community}

\subsection{\textcolor{red}{(Geometric Kernel Block Model, GKBM)}}
\citep{avrachenkov2024community} 
\citep{allem2025multicommunityspectralclusteringgeometric}

\subsection{\textcolor{red}{Geometric Hidden Community Model (GHCM)}}

\citep{gaudio2024exactcommunityrecoverygeometric} \citep{gaudio2024exact}
\citep{gaudio2025sharpexactrecoverythreshold}
\citep{gaudio2026exact} 

\section{\textcolor{red}{Spectral}}
(A confirmer) \citep{chen2016community}

(papier SBM mais théorie intéressante) \citep{bordenave2015non}

Cadre spectral intéressant pour certains espaces \citep{meliot2019asymptotic}

Méthode spectrale + Robustesse (la perturbation est géométrique) \citep{peche2020robustness}

Méthode spectrale sur Soft Geometric Block Model (SGBM) (Generalisation de Geomemetric SBM)
\citep{avrachenkov2021higher}



\section{Ouvertures}
\begin{itemize}
    \item Considérer des modèles hyperboliques, pas seulement labels \& distance-dependant,
    car ceux-ci captent mieux les propriétés habituelles des réseaux 
    \begin{itemize}
        \item Clustering \citep{Krioukov_2016}
        \item Degre inhomogène \citep{Baraba_si_1999}
        \item "Communities with specific mixing patterns" (mesoscale) \citep{Toivonen_2006} 
    \end{itemize} 
    Communautés apparaissent dans un espace hyperbolique latent \citep{bruno2019community}
    Réponse 6ans plus tard : pas toujours ! + definition d'un modele pour "régler" ça
    \citep{guarino2025random}
% Bruno et al. : 
% “les communautés peuvent émerger dans un embedding hyperbolique appris” ;
% Guarino et al. :
% “oui, mais la géométrie hyperbolique pure ne suffit pas toujours à expliquer la
% structure mésoscopique ; il faut parfois ajouter une vraie composante de blocs.”

\citep{Li_Sha_Huang_Zhang_2018}
% au-delà des graphes géométriques, d’autres travaux incorporent
% des informations latentes non topologiques, par exemple
% des attributs de nœuds, souvent via des approches
% d'embedding ou de factorisation non négative



\end{itemize}

% Springer author-year bibliography style
\bibliographystyle{plainnat}
\bibliography{../pdf/refs_filtered}

\end{document}
